% ===================================================================
%  नकसा नं. १६ — बड़का दुकान। --- बजार। खेलौना।
%
%  Traduction, faite en 2026, en bhojpouri d'aujourd'hui, sur l'ido
%  de texto/io. Regles de la colonne : voir 10-nakasa-01.tex. Une
%  seule numerotation. 33 blocs, 106 renvois, dix-huit paires a
%  retourner. DERNIER TABLEAU DU LIVRET, ET DERNIERE DES
%  VINGT-NEUF LANGUES. UNE note d'auteur, t16-noto-1, posee entre les
%  blocs 8 et 9 ; son appel se compose en parentheses simples, « (*) ».
%  L'ido coupe cinq alineas sur des « suite » de feuillet ; la
%  traduction n'herite d'aucune des cinq.
%
%  ECART DECLARE : t16-15-1. L'ido saute la lettre « (d) » que le
%  fac-simile grave pourtant ; renvoji.py laisse donc ce bloc passer
%  SANS RIEN VERIFIER. L'ordre a ete compte a la main : 85, a, b, c,
%  d, e, f, 86, 87, 88, 89 — onze appels.
%
%  LE TITRE N'A PAS BESOIN D'ETRE TRADUIT, ET C'EST LA QUATRIEME FOIS
%  QUE LE RELEVE LE CONSTATE AU MEME ENDROIT. « La Bazaro » s'ecrit
%  ici « बजार » ; la colonne gujaratie ecrivait « બજાર » et la colonne
%  levantine « البازار ». Un seul mot persan, bāzār, parti de Chiraz
%  vers l'est par les Ghurides et les Moghols, vers l'ouest par les
%  Seldjoukides puis l'italien ; et le francais de Rochelle ecrit
%  « bazar » sans s'apercevoir qu'il ecrit du persan. TROIS COLONNES
%  ET L'ORIGINAL SUR UN SEUL MOT, AU DERNIER TITRE DU LIVRET, POUR LA
%  DERNIERE DES VINGT-NEUF LANGUES.
%
%  L'HIPPOPOTAME (22) EST « दरियाई घोड़ा », ET C'EST LE MEME CALQUE
%  QUE LA COLONNE LEVANTINE. Elle avait releve « فرس النهر », le
%  cheval du fleuve, calque mot pour mot du grec ἱπποπόταμος par les
%  traducteurs de Bagdad. Le bhojpouri dit dariyā, le fleuve, et
%  ghoṛā, le cheval : la meme image, arrivee par une autre route et
%  dans une autre famille de mots. DEUX LANGUES QUI N'ONT JAMAIS VU
%  LA BETE ONT INVENTE POUR ELLE LA MEME METAPHORE. A cote, le
%  rhinoceros (7) n'a pas eu besoin de metaphore : गैंडा est du vieux
%  fonds, parce que l'animal vit dans le terai, a trois cents
%  kilometres de Chapra.
%
%  ET LE TABLEAU D'ASTRONOMIE (85) EST LE SEUL ENDROIT DU LIVRET OU
%  CETTE LANGUE N'EMPRUNTE RIEN A PERSONNE. Le tableau 6 avait montre
%  la lumiere nommee par l'anglais et le persan, le tableau 9 la
%  montagne nommee par le sanscrit savant, le tableau 12 le chemin de
%  fer nomme par l'anglais d'un bout a l'autre. Ici : धूमकेतु la
%  comete, ग्रहण l'eclipse, गोलार्ध l'hemisphere, ग्रह les planetes,
%  सौरमंडल le systeme solaire — du sanscrit, mais du sanscrit qui n'a
%  jamais rien traduit, parce que l'astronomie indienne est plus
%  vieille que celle qu'on enseigne sur cette planche. Et les quatre
%  points cardinaux (c) — उत्तर, दक्खिन, पूरब, पच्छिम — sont du fonds
%  le plus ancien : le tableau 8 avait ecrit पछुआ pour le vent
%  d'ouest, et c'est la meme racine. LA LANGUE EMPRUNTE POUR LES
%  CHOSES ET NON POUR LE CIEL.
% ===================================================================

%%K t16-apar-1 apar
\VUsaut{4.0mm}
\VUtitre{15.0pt}{60}{नकसा नं. १६}
\VUsaut{2.0mm}
\VUpk{13.2pt}{40}{बड़का दुकान। --- बजार।}
\VUpk{13.2pt}{40}{खेलौना।}
\VUsaut{3.4mm}

%%K t16-01-1 p
१. --- नया साल नगीच आ रहल बा। बड़का दुकान आपन बसता \VUgras{खेलौना}
\textsuperscript{(1)} आ नया साल के \VUgras{भेंट} से सजा लेले बा।

%%K t16-01-2 p
हम आपन बाबा से कहनी कि हमरा के ई सुंदर नुमाइस देखे बजार ले चलीं, आ ऊ
मान गइलें।

%%K t16-02-1 p
२. --- पहिले हमनी के सुंदर हथियार के तखता के सामने रुकनी जा, जवना में
\VUgras{बंदूक}, \VUgras{फौजी टोपी}, \VUgras{तलवार}, \VUgras{तलवार के
पेटी} आ एक जोड़ी \VUgras{कान्ह के पट्टी} रहे, भा \VUgras{खोद},
\VUgras{कवच} \textsuperscript{(3)} आ \VUgras{पिस्तौल} \textsuperscript{(4)}।
ई सब हमरा के \VUgras{जादुई लालटेन} \textsuperscript{(2)} से बहुते जादे नीक
लागल।

%%K t16-tit-1 sub
\VUsaut{3.0mm}
\VUpk{10.2pt}{40}{सुंदर नजारा।}
\VUsaut{2.6mm}

%%K t16-03-1 p
३. --- ओही विभाग में एगो \VUgras{पहरेदार} \textsuperscript{(5)} आपन
\VUgras{पहरा के गुमटी} में रहे ; लागत रहे कि ऊ कागज के पूरा जनावरखाना
के पहरा देत बा। ओहिजा केतना जनावर रहे ! एगो \VUgras{तोता}
\textsuperscript{(6)} आपन अड्डा पर, एगो \VUgras{गैंडा}
\textsuperscript{(7)} नाक पर सींग लेके, एगो \VUgras{रेनडियर}
\textsuperscript{(8)}, एगो \VUgras{सुतुरमुर्ग}
\textsuperscript{(9)}, एगो \VUgras{बंदर} \textsuperscript{(10)} जे अखरोट
तोड़त बा, एगो \VUgras{लोमड़ी} \textsuperscript{(11)} जे मुरगी ले भागत बा,
एगो \VUgras{मगर} \textsuperscript{(12)} जे आपन बड़हन \VUgras{जबड़ा} फारले
बा, एगो \VUgras{ऊँट} \textsuperscript{(13)}, एगो सुघर \VUgras{सलूकी}
\textsuperscript{(14)}, एगो \VUgras{तेंदुआ} \textsuperscript{(15)}, आ ओकरा
बाद दू गो जनावर जेकरा के हम ना जानत रहनी आ जे, लोग के कहला के मुताबिक,
\VUgras{वीजल} \textsuperscript{(16)} आ \VUgras{लकड़बग्घा} \textsuperscript{(17)}
हउवें। ओहिजा एगो \VUgras{चीता} \textsuperscript{(18)} आ एगो
\VUgras{भेड़िया} \textsuperscript{(21)} भी रहे ; एकदम लगे रबर के सुघर छोट
\VUgras{चूहा} \textsuperscript{(19)} आ बहुते छोट \VUgras{मूस}
\textsuperscript{(20)} रहे। आखिर में एगो \VUgras{दरियाई घोड़ा}
\textsuperscript{(22)} आ एगो कूबड़ वाला \VUgras{एक-कूबड़ ऊँट}
\textsuperscript{(23)} एह मजमा के पूरा करत रहे। हम ओहिजा आउर कई घंटा रह
जइतीं जो लगहीं \VUgras{सीसा के सिपाही} \textsuperscript{(29)} से भरल एगो
शानदार पड़ाव हमार धेयान ना खींच लेत।

%%K t16-tit-2 sub
\VUsaut{4.0mm}
\VUpk{10.2pt}{40}{सिपाही आ लड़ाई।}
\VUsaut{2.8mm}

%%K t16-04-1 p
४. --- हमार बाबा, जे \VUgras{अपाहिज फौजी} \textsuperscript{(24)} हउवें आ
जेकरा एगो \VUgras{घाव} के चलते फौज छोड़े के पड़ल, जवना से ओकरा
\VUgras{फौजी तमगा} मिलल, हमरा के ओह \VUgras{लड़ाई} के हाल सुनावल बहुते
पसंद करेलें जवना में ऊ लड़ल रहलें। ऊ खुसी से हमरा के ओह छोट मिनिएचर
फौज के अलग-अलग चाल समझावत रहलें। बजार देखे आइल सच्चा सिपाही लोग के एगो
टोली — एगो \VUgras{इंजीनियर सिपाही} \textsuperscript{(25)}, एगो
\VUgras{पहाड़ी सिकारी} \textsuperscript{(26)} जे \VUgras{बेरे टोपी}
\textsuperscript{(27)} पहिनले रहे, पैदल सेना के एगो \VUgras{हवलदार-मेजर}
\textsuperscript{(28)} आ तोपखाना के एगो \VUgras{हवलदार} आपन कमीज पर सोना के
\VUgras{फीता} लगववले — हमनी के लगे रुक के ई समझावल सुने लागल।

%%K t16-05-1 p
५. --- हमार बाबा, एतना नीक सुनइया पाके बहुते फूल गइलें, आ फउरन एगो
पूरा लड़ाई के नकसा बना देलें।

%%K t16-05-2 p
आपन \VUgras{परकोटा} आ \VUgras{खाई} के साथे किलाबंद सहर के \VUgras{दुसमन
के फौज} घेरले रहे, आ लमहर \VUgras{गोलाबारी} के बाद ऊ \VUgras{धावा} बोले
के तइयार रहे। बाकिर \VUgras{घिराइल} लोग, भूख से तंग आके, \VUgras{घेरे
वाला} के \VUgras{पड़ाव} पर \VUgras{उलटा धावा} बोले के कोसिस कइलस।
\VUgras{तंबू} के चारो ओर के जिद्दी \VUgras{लड़ाई} में \VUgras{हमला}
पीछे धकेलला के बाद \VUgras{जितइया} घेरे वाला आपन \VUgras{घोड़सवार दस्ता}
के \VUgras{हारल} हमलावर लोग के \VUgras{पीछा} पर लगा देलस। ई लोग
\VUgras{पीछे हटल}, आ \VUgras{रनभूमि} के \VUgras{मुरदा} आ \VUgras{घायल}
से भर देलस। \VUgras{नक्कालात वाला} लोग घायलन के \VUgras{एंबुलेंस} में
ले गइल ताकि \VUgras{जर्राह} ओकनी के इलाज करसु।

%%K t16-05-3 p
कुछ \VUgras{दस्ता} के बहादुर बिरोध के बादो, जे \VUgras{चौकोर} बना के
ठाढ़ रहे, \VUgras{हार} पूरा भइल, बाकी फौज \VUgras{भाग} गइल आ बहुते
\VUgras{बंदी} छोड़ गइल। दुसमन सहर पर कब्जा करके आपन \VUgras{जीत} पूरा
करे वाला रहे। सायद इहे लड़ाई के अंत होई आ \VUgras{हथियार बंदी} के बाद
\VUgras{सांति के संधि} पर दसखत हो सकी।

%%K t16-tit-3 sub
\VUsaut{4.4mm}
\VUpk{10.2pt}{40}{नया साल के भेंट।}
\VUsaut{2.8mm}

%%K t16-06-1 p
६. --- जवान सिपाही लोग, जे ओह पुरनका लड़इया के रंगीन कहानी सुन के
हँसत रहे, ओकरा से कुछ आउर बात पूछलस। हम त दुकान के चारो ओर घूम के एगो
सुंदर \VUgras{क्रॉसबो} \textsuperscript{(32)} आ एगो \VUgras{धनुष}
\textsuperscript{(33)} आपन \VUgras{तीर} \textsuperscript{(31)} के साथे देखत
रहनी। ओहिजा हमरा जनावरन के एगो दोसर मजमा मिलल : दू गो \VUgras{गवइया
चिरई} \textsuperscript{(34)} — एगो मसीनी \VUgras{बुलबुल} आ एगो
\VUgras{फुदकी} जे गला फार के गावत रहे — आ अजीब जनावरन के एगो कतार :
\VUgras{चमगादड़} \textsuperscript{(35)}, \VUgras{अजगर} \textsuperscript{(36)},
\VUgras{कछुआ} \textsuperscript{(37)}, \VUgras{सील} \textsuperscript{(38)},
\VUgras{व्हेल} \textsuperscript{(39)} आ फुलावल रबर के \VUgras{शार्क}
\textsuperscript{(40)}।

%%K t16-07-1 p
७. --- हम ओकनी के देखे में देर ना कइनी, काहेकि हमरा ऊ लउक गइल जेकर हम
सपना देखत रहनी : एगो शानदार \VUgras{कठपुतली के नाटकघर} \textsuperscript{(41)}
आपन बढ़िया \VUgras{साज-सज्जा} आ मन मोह लेवे वाला लाल \VUgras{परदा} के
साथे। \VUgras{मंच} पर दू गो \VUgras{नट} एगो बहुते मजेदार \VUgras{नाटक}
में आपन \VUgras{भूमिका} निभावत लागत रहलें, काहेकि कागज के छोट
\VUgras{दरसक}, जे \VUgras{मचान} में बइठल रहे भा \VUgras{कुरसी} पर आ
\VUgras{बजनियाँ के मुखिया} के पाछे \VUgras{अगला जगहा} में, जोर से
ताली बजावत रहे।

%%K t16-07-2 p
अफसोस, ई नाटकघर बहुते महँग रहे।

%%K t16-08-1 p
८. --- ऊपर से, खेलौना चुनल मुसकिल रहे, काहेकि सो-केस में हर तरह के
खेलौना जमा रहे : \VUgras{अर्लेकिन} \textsuperscript{(42)},
\VUgras{पोलिशिनेल} \textsuperscript{(43)}, \VUgras{पिएरो}
\textsuperscript{(44)}, \VUgras{उल्लू} \textsuperscript{(45)} जे पाँख
फड़फड़ावत रहे, सीटी बजावे वाला \VUgras{कस्तूरा} \textsuperscript{(46)},
भूँके वाला \VUgras{पूडल कुकुर}, एगो \VUgras{फोनोग्राफ} \textsuperscript{(47)}
आ \VUgras{धीरज के खेल}। ओहमें से एगो खेल हमरा बहुते नीक लागल : ओहमें एगो
\VUgras{समुंदरी लड़ाई} \textsuperscript{(48)} देखावल गइल रहे जवना में एगो
\VUgras{जहाज} पर \VUgras{समुंदरी लुटेरा} लोग हमला करत रहे, एगो अइसन
जमीन लगे जहाँ \VUgras{नरियर के गाछ} लागल रहे।

%%K t16-noto-1 noto
\VUnotes{143.2mm}{%
(*) नकसा नं. ६ देखीं।%
}

%%K t16-09-1 p
९. --- हम ई सब अचरज बहुते आराम से देख पवनी, काहेकि दुकान में थोड़े लोग
रहे — बस कुछ टहले वाला, एगो \VUgras{घोड़सवार सिपाही} \textsuperscript{(49)}
आ एगो \VUgras{वसूली करे वाला} \textsuperscript{(50)} जे मुख्य \VUgras{गल्ला}
\textsuperscript{(51)} पर एगो \VUgras{हुंडी} देत रहे।

%%K t16-10-1 p
१०. --- हम आपन बाबा से पूछनी कि गल्ला वाली के पाछे तखता पर धइल किताब
कवना काम के हउवें ; ऊ जवाब देलें कि ऊ \VUgras{बही-खाता} हउवें।

%%K t16-tit-4 sub
\VUsaut{3.6mm}
\VUpk{10.2pt}{40}{दुकान के सामने ताकल।}
\VUsaut{2.8mm}

%%K t16-11-1 p
११. --- खेलौना के विभाग छोड़ के हमनी के जीन के विभाग में गइनी जा ताकि
\VUgras{जीन} \textsuperscript{(53)}, \VUgras{रकाब} \textsuperscript{(54)},
\VUgras{लगाम} \textsuperscript{(55)}, \VUgras{दहाना} \textsuperscript{(56)}
आ \VUgras{चाबुक} देख सकीं जा। जब \VUgras{विभाग के मुखिया}
\textsuperscript{(57)} हमार बाबा के कुछ बात बतावत रहलें, हम \VUgras{चीनी
मिट्टी आ काँच} \textsuperscript{(58)} के बसता देखे चल गइनी जहाँ सुंदर
\VUgras{सलाद के कटोरा} आ नाजुक \VUgras{चाय के केतली} \textsuperscript{(59)}
रहे।

%%K t16-12-1 p
१२. --- पहिला मंजिल पर जाए खातिर हमनी के \VUgras{कोरसेट}
\textsuperscript{(60)} के सो-केस के पाछे से गुजरनी जा, मोजा के दुकान लगे,
जहाँ बहुते सुंदर \VUgras{जाँघिया} \textsuperscript{(63)} सजल रहे। लगहीं,
एगो \VUgras{दुकान के लइकी} \textsuperscript{(61)} एगो \VUgras{खरीदार
मेहरारू} \textsuperscript{(62)} के सुंदर रेसमी \VUgras{कपड़ा}
\textsuperscript{(64)} चुनवावत रहली।

%%K t16-12-2 p
सीढ़ी से चढ़त हम देखनी कि दुकान जापानी \VUgras{लालटेन} \textsuperscript{(65)},
\VUgras{छाता} \textsuperscript{(66)} आ बहुते बड़का \VUgras{खजूर के गाछ}
\textsuperscript{(70)} से सजल बा।

%%K t16-13-1 p
१३. --- पहिला मंजिल पर देखे लायक कुछ जादे ना रहे ; बस फर्नीचर (*),
\VUgras{लोहा के तिजोरी} \textsuperscript{(67)}, \VUgras{दराज वाला आलमारी}
\textsuperscript{(68)} आ \VUgras{बच्चा के गाड़ी} \textsuperscript{(69)}।
बाकिर पूरा दुकान के नजारा बहुते \VUgras{मजेदार} रहे।

%%K t16-14-1 p
१४. --- हमनी के गोड़ के नीचे हम बाजा देखनी : \VUgras{वीणा}
\textsuperscript{(71)}, \VUgras{गिटार} \textsuperscript{(72)},
\VUgras{मैंडोलिन} \textsuperscript{(73)}, \VUgras{पिस्टन वाला भोंपू}
\textsuperscript{(74)}, \VUgras{क्लैरिनेट} \textsuperscript{(75)}, आपन
\VUgras{गज} \textsuperscript{(76)} के साथे बेला आ इहाँ तक कि एगो
\VUgras{बड़का ढोल} \textsuperscript{(77)}। थोड़ा आगे, कागज के विभाग लगे,
एगो \VUgras{बिद्यार्थी} \textsuperscript{(78)} एगो \VUgras{दुकान के छोकड़ा}
\textsuperscript{(79)} से कापी के थैली के भाव-ताव करत रहे। ओह लोग लगे हम
एगो सुंदर \VUgras{वर्णमाला के किताब} \textsuperscript{(80)} देखनी।

%%K t16-14-2 p
दुकान के बीचोबीच एगो ऊँच \VUgras{बड़ दिन के गाछ} \textsuperscript{(81)}
खड़ा रहे, जवन मोमबत्ती, छोट सजावट आ हर तरह के खेलौना से भरल रहे :
\VUgras{काठ के घोड़ा} \textsuperscript{(82)}, \VUgras{गुड़िया}
\textsuperscript{(83)}, वगैरह।

%%K t16-14-3 p
\VUgras{इतर के दुकान} \textsuperscript{(84)} आपन \VUgras{दाँत के पानी} आ
तरह-तरह के \VUgras{इतर} के साथे बहुते नीक महक फइलावत रहे जवन हमनी तक
पहुँचत रहे।

%%K t16-tit-5 sub
\VUsaut{3.4mm}
\VUpk{10.2pt}{40}{हमनी के कुछ खरीदत बानी जा।}
\VUsaut{2.8mm}

%%K t16-15-1 p
१५. --- काहेकि हमार बाबा के लागे लागल कि बहुते देर हो गइल, हमनी के
बड़का सीढ़ी से नीचे उतरनी जा। ऊ एगो \VUgras{खगोल के तखता}
\textsuperscript{(85)} के सामने थोड़ा रुकलें, जवन — ऊ हमरा के कहलें —
\VUgras{धूमकेतु} \textsuperscript{(a)}, \VUgras{चंद्रगहन आ सूरजगहन}
\textsuperscript{(b)}, \VUgras{चार दिसा} \textsuperscript{(c)}
\VUgras{(उत्तर, दक्खिन, पूरब आ पच्छिम)} वाला हवा-चक्र, दू गो
\VUgras{गोलार्ध} \textsuperscript{(d)} के नकसा, \VUgras{ग्रह}
\textsuperscript{(e)} आ \VUgras{सौरमंडल} \textsuperscript{(f)} देखावत बा।
एह सब में हमरा कुछ समझ ना आइल, आ हमरा त गुजरत \VUgras{पहुँचावे वाला
छोकड़ा} \textsuperscript{(86)} के फीता वाला वरदी में बहुते जादे मन लागल,
आ ओह सुंदर \VUgras{पंखा} \textsuperscript{(87)} में जे हमरा लगे रहे,
\VUgras{बटुआ} \textsuperscript{(88)} आ \VUgras{पर्स} \textsuperscript{(89)}
के बगल में।

%%K t16-16-1 p
१६. --- हम बसता पर \VUgras{पंख} \textsuperscript{(90)} आ \VUgras{पेटी के
बकुआ} \textsuperscript{(91)} भी निहारनी, आ ऊ सुंदर \VUgras{बनावटी फूल}
\textsuperscript{(94)} जे टोकरी में सुघर ढंग से सजल रहे। \VUgras{बनफसा},
\VUgras{तुदरी}, \VUgras{सुंबुल} \textsuperscript{(95)},
\VUgras{जिरेनियम} \textsuperscript{(96)}, \VUgras{लिली} \textsuperscript{(97)}
आ \VUgras{नस्तूर्तियम} \textsuperscript{(98)} बहुते नीक से नकल कइल गइल
रहे।

%%K t16-tit-6 sub
\VUsaut{4.4mm}
\VUpk{10.2pt}{40}{बाबा के भेंट।}
\VUsaut{2.8mm}

%%K t16-17-1 p
१७. --- हम आपन बाबा से बिनती कइनी कि हमरा खातिर ओह \VUgras{दाँत के
बुरुस} \textsuperscript{(92)} में से एगो खरीद दीं जे \VUgras{बाल के पिन}
\textsuperscript{(93)} लगे एगो डिब्बा में रहे। ऊ मान गइलें आ हम खुदे आपन
सामान के दाम गल्ला पर देवे गइनी, जेकरा लगे एगो बड़हन \VUgras{खेप}
\textsuperscript{(100)} एगो \VUgras{ग्राहक} \textsuperscript{(99)} खातिर
तइयार कइल जात रहे।

%%K t16-18-1 p
१८. --- अचानक ओह लोग के बीच थोड़ा हलचल भइल जे कलात्मक \VUgras{काँसा के
मूरत} \textsuperscript{(101)} लगे रुकल रहे। एगो \VUgras{चोर}
\textsuperscript{(102)} के अभिए एगो \VUgras{निरीक्षक} \textsuperscript{(103)}
रंगे हाथ धर लेले रहलें, जब ऊ केहू के जेब काटे के कोसिस करत रहे। लोग
ओकर \VUgras{गिरफ्तारी} खातिर पुलिस बोलववलस, आ ठीक जब हमनी के बाहर
निकलत रहनी जा, ओह गुनहगार के हवालात ले जात रहे।
